\xchapter{To Flavian commonly called “the Tome.”}

Letter XXVIII.

I. \emph{Eutyches has been driven into his error by presumption and ignorance}\footnote{The original word (\emph{imperitia}) implies that a recluse like Eutyches (an archimandrite of a convent) ought never to have entered into a nice controversy like the present: he has not enough \emph{savoir faire}, and his knowledge is not quite up to date, is a little old-fashioned.}.

Having read your letter, beloved, at the late arrival of which we are surprised\footnote{The exact reason of the delay is not altogether certain: we know Flavian had written much earlier than the date of arrival warranted: it is No. XXII. in the series.}, and having perused the detailed account of the bishops’ acts\footnote{Viz., the proceedings of the \textgreek{σύνοδος ἐνδημοῦσα} summoned by Flavian at Constantinople.}, we have at last found out what the scandal was which had arisen among you against the purity of the Faith: and what before seemed concealed has now been unlocked and laid open to our view: from which it is shown that Eutyches, who used to seem worthy of all respect in virtue of his priestly office, is very unwary and exceedingly ignorant, so that it is even of him that the prophet has said: “he refused to understand so as to do well: he thought upon iniquity in his bed\footnote{Ps. xxxvi. 4.}.” But what more iniquitous than to hold blasphemous opinions\footnote{\emph{Impia sapere}, to think disloyal things against God: cf. the \emph{recta sapere}, “to have a right judgment” of the Collect for Whitsunday.}, and not to give way to those who are wiser and more learned than ourself. Now into this unwisdom fall they who, finding themselves hindered from knowing the truth by some obscurity, have recourse not to the prophets’ utterances, not to the Apostles’ letters, nor to the injunctions of the Gospel but to their own selves: and thus they stand out as masters of error because they were never disciples of truth. For what learning has he acquired about the pages of the New and Old Testament, who has not even grasped the rudiments of the Creed? And that which, throughout the world, is professed by the mouth of every one who is to be born again\footnote{Knowledge of and belief in the principles of the Faith as contained in the Creed (\emph{symbolum}) have of course always been required before Baptism from very early times. Leo here calls catechumens \emph{regenerandi}, just as those who are being baptized are spoken of as \emph{renascentes} (e.g. Lett. XVII. 8), those who have been baptized as \emph{renati} (\emph{passim}), and the rite itself as \emph{sacramentum regenerationis} (e.g. Lett. IX. 2).}, is not yet taken in by the heart of this old man.

II. \emph{Concerning the twofold nativity and nature of Christ.}

Not knowing, therefore, what he was bound to think concerning the incarnation of the Word of \textsc{God}, and not wishing to gain the light of knowledge by researches through the length and breadth of the Holy Scriptures, he might at least have listened attentively to that general and uniform confession, whereby the whole body of the faithful confess that they \emph{believe in \textsc{God} the Father Almighty, and in Jesus Christ, His only Son}\footnote{The Latin \emph{unicus} is not so exact as the Greek original \textgreek{μονογενής}: elsewhere, however, \emph{unigenitus} is used.}\emph{, our \textsc{Lord}, who was born of the Holy Spirit and}\footnote{N.B. \emph{et} (and) not \emph{ex} (out of).}\emph{the Virgin Mary.} By which three statements the devices of almost all heretics are overthrown. For not only is \textsc{God} believed to be both Almighty and the Father, but the Son is shown to be co-eternal with Him, differing in nothing from the Father because He is \emph{\textsc{God} from \textsc{God}}\footnote{The language of the Nicene Creed.}, Almighty from Almighty, and being born from the Eternal one is co-eternal with Him; not later in point of time, not lower in power, not unlike in glory, not divided in essence: but at the same time the only begotten of the eternal Father was born eternal of the Holy Spirit and the Virgin Mary. And this nativity which took place in time took nothing from, and added nothing to that divine and eternal birth, but expended itself wholly on the restoration of man who had been deceived\footnote{I.e. by the Devil: the allusion is to Adam’s fall in Paradise.}: in order that he might both vanquish death and overthrow by his strength\footnote{\emph{Sua virtute}: in patristic Latin \emph{virtus} is, as is well known, usually the translation of the Greek \textgreek{δύναμις} and has a much wider meaning than moral excellence, our virtue.}, the Devil who possessed the power of death. For we should not now be able to overcome the author of sin and death unless He took our nature on Him and made it His own, whom neither sin could pollute nor death retain. Doubtless then, He was conceived of the Holy Spirit within the womb of His Virgin Mother, who brought Him forth without the loss of her virginity, even as she conceived Him without its loss.

But if he could not draw a rightful understanding (of the matter) from this pure source of the Christian belief, because he had darkened the brightness of the clear truth by a veil of blindness peculiar to himself, he might have submitted himself to the teaching of the Gospels. And when Matthew speaks of “the Book of the Generation of Jesus Christ, the Son of David, the Son of Abraham\footnote{S. Matt. i. 1.},” he might have also sought out the instruction afforded by the statements of the Apostles. And reading in the Epistle to the Romans, “Paul, a servant of Jesus Christ, called an Apostle, separated unto the Gospel of \textsc{God}, which He had promised before by His prophets in the Holy Scripture concerning His son, who was made unto Him\footnote{\emph{ei}. So the Vulgate.} of the seed of David after the flesh\footnote{Rom. i. 1–3.},” he might have bestowed a loyal carefulness upon the pages of the prophets. And finding the promise of \textsc{God} who says to Abraham, “In thy seed shall all nations be blest\footnote{Gen. xii. 3.},” to avoid all doubt as to the reference of this seed, he might have followed the Apostle when He says, “To Abraham were the promises made and to his seed. He saith not and to seeds, as if in many, but as it in one, and to thy seed which is Christ\footnote{Gal. iii. 16.}.” Isaiah’s prophecy also he might have grasped by a closer attention to what he says, “Behold, a virgin shall conceive and bear a Son and they shall call His name Immanuel,” which is interpreted “\textsc{God} with us\footnote{Is. vii. 14. and S. Matt. i. 23.}.” And the same prophet’s words he might have read faithfully. “A child is born to us, a Son is given to us, whose power is upon His shoulder, and they shall call His name the Angel of the Great Counsel, Wonderful, Counsellor, the Mighty \textsc{God}, the Prince of Peace, the Father of the age to come\footnote{Is. ix. 6. “The angel of the great counsel” (\emph{magni consilii angelus}) is a translation of the LXX. (which in the rest of the verse either represents a very different original text, or contents itself with a loose paraphrase), and is again repeated in the “Counsellor” (\emph{Consiliarius}), two words farther on (which is also the Vulgate reading).}.” And then he would not speak so erroneously as to say that the Word became flesh in such a way that Christ, born of the Virgin’s womb, had the form of man, but had not the reality of His mother’s body\footnote{This was the third dogma of Apollinaris (more fully stated in Lett. CXXIV. 2 and CLXV. 2) that our Lord’s acts and sufferings as man belonged entirely to His Divine nature, and were not really human at all.}. Or is it possible that he thought our \textsc{Lord} Jesus Christ was not of our nature for this reason, that the angel, who was sent to the blessed Mary ever Virgin, says, “The Holy Ghost shall come upon thee and the power of the Most High shall overshadow thee: and therefore that Holy Thing also that shall be born of thee shall be called the Son of \textsc{God\footnote{S. Luke i. 35.}},” on the supposition that as the conception of the Virgin was a Divine act, the flesh of the conceived did not partake of the conceiver’s nature? But that birth so uniquely wondrous and so wondrously unique, is not to be understood in such wise that the properties of His kind were removed through the novelty of His creation. For though the Holy Spirit imparted fertility to the Virgin, yet a real body was received from her body; and, “Wisdom building her a house\footnote{Prov. ix. 1.},” “the Word became flesh and dwelt in us\footnote{\emph{In nobis}, which he seems from the immediately following words to interpret as meaning “in our flesh,” and not “amongst us,” as the R.V. and others.},” that is, in that flesh which he took from man and which he quickened with the breath of a higher life\footnote{\emph{Quam spiritu vitæ rationalis} (\textgreek{λογικοῦ}) \emph{animavit}.}.

III. \emph{The Faith and counsel of \textsc{God} in regard to the incarnation of the Word are set forth.}

Without detriment therefore to the properties of either nature and substance which then came together in one person\footnote{A famous passage quoted by Hooker, Eccl. Pol. v. 53, 2, and Liddon Bampt. Lect., p. 267. Compare Serm. lxii. 1, \emph{quod…in unam personam concurrat proprietas utriusque substantiæ} (Bright), also xxii. 2, xxiii. 2.}, majesty took on humility, strength weakness, eternity mortality: and for the paying off of the debt belonging to our condition inviolable nature was united with possible nature, so that, as suited the needs of our case\footnote{\emph{Quod nostris remediis congruebat}, where \emph{remedia} must mean the disease which needs remedies (a sort of passive use).}, one and the same Mediator between \textsc{God} and men, the Man Christ Jesus, could both die with the one and not die with the other. \footnote{This passage from “Thus in the whole” to “not the failing of power” is repeated again in Sermon xxiii. 2, almost word for word.}Thus in the whole and perfect nature of true man was true \textsc{God} born, complete in what was His own, complete in what was ours. And by “ours” we mean what the Creator formed in us from the beginning and what He undertook to repair. For what the Deceiver brought in and man deceived committed, had no trace in the Saviour. Nor, because He partook of man’s weaknesses, did He therefore share our faults. He took the form of a slave\footnote{The reference, of course, is to Phil. ii. 6: no passage is a greater favourite with the Fathers than this.} without stain of sin, increasing the human and not diminishing the divine: because that emptying of Himself whereby the Invisible made Himself visible and, Creator and \textsc{Lord} of all things though He be, wished to be a mortal, was the bending down\footnote{Compare S. Aug. ad Catech. § 6, \emph{humilitas Christi quid est? manum Deus homini iacenti porrexit: nos cecidimus, ille descendit: nos iacebamus, ille se inclinavit. Prendamus et surgamus ut non in pœnam cadamus.}} of pity, not the failing of power. Accordingly He who while remaining in the form of \textsc{God} made man, was also made man in the form of a slave. For both natures retain their own proper character without loss: and as the form of \textsc{God} did not do away with the form of a slave, so the form of a slave did not impair the form of \textsc{God}. For inasmuch as the Devil used to boast that man had been cheated by his guile into losing the divine gifts, and bereft of the boon of immortality had undergone sentence of death, and that he had found some solace in his troubles from having a partner in delinquency\footnote{\emph{De prævaricatoris consortio}: \emph{prævaricator} originally is a legal term, signifying “a shuffler” in a suit, an advocate who plays into the hands of the other side.}, and that \textsc{God} also at the demand of the principle of justice had changed His own purpose towards man whom He had created in such honour: there was need for the issue of a secret counsel, that the unchangeable \textsc{God} whose will cannot be robbed of its own kindness, might carry out the first design of His Fatherly care\footnote{\emph{Pietas}, as in the collect for xvi. S. aft. Trin., where the English, “pity” represents the Latin “\emph{pietas}” philologically as well as in meaning. Cf. n. 2 in chap. vi.} towards us by a more hidden mystery\footnote{\emph{Sacramento}, (\textgreek{μυστηρίῳ}): what the “mystery” was is finely set forth by Canon Bright’s hymn, No. 172, H. A. and M. (new edition).}; and that man who had been driven into his fault by the treacherous cunning of the devil might not perish contrary to the purpose of \textsc{God\footnote{The whole of the end of this chapter from “For inasmuch as,” and the beginning of the next down to “laws of death,” is repeated word for word in Sermon XXII., chaps. i. and ii.}}.

IV. \emph{The properties of the twofold nativity and nature of Christ are weighed one against another.}

There enters then these lower parts of the world the Son of \textsc{God}, descending from His heavenly home and yet not quitting His Father’s glory, begotten in a new order by a new nativity. In a new order, because being invisible in His own nature, He became visible in ours, and He whom nothing could contain was content to be contained\footnote{\emph{Incomprehensibilis voluit comprehendi}. Canon Bright’s references are most apposite: “compare Serm. lxviii., \emph{idem est qui impiorum manibus comprehenditur et qui nullo fine concluditur:} and Serm. xxxvii. 1, \emph{genetricis gremio continetur qui nullo fine concluditor}. This ‘antithesis’ has been grandly expressed in Milman’s, ‘Martyr of Antioch.’ \par “‘And Thou wast laid within the tomb… \par Whom heaven could not contain, \par Nor the immeasurable plain \par Of vast infinity enclose or circle round.’”}: abiding before all time He began to be in time: the \textsc{Lord} of all things, He obscured His immeasurable majesty and took on Him the form of a servant: being \textsc{God} that cannot suffer, He did not disdain to be man that can, and, immortal as He is, to subject Himself to the laws of death. The \textsc{Lord} assumed His mother’s nature without her faultiness: nor in the \textsc{Lord} Jesus Christ, born of the Virgin’s womb, does the wonderfulness of His birth make His nature unlike ours. For He who is true \textsc{God} is also true man: and in this union there is no lie\footnote{I.e. , there is no fancy, no pretending: each nature is in equal reality present, the human as well as the Divine, thus opposing all Docetic and Monophysite heresies.}, since the humility of manhood and the loftiness of the Godhead both meet there. For as \textsc{God} is not changed by the showing of pity, so man is not swallowed up by the dignity. For each form does what is proper to it with the co-operation of the other\footnote{This passage (which is repeated in Serm. liv., chap. 2, down to “injuries”), was objected to by the Illyrian and Palestinian bishops as savouring of the heresy of Nestorius who “divided the substance:” but it is obvious that the same words might have an orthodox meaning in the mouth of one who was orthodox and to the unorthodox would bear an unorthodox construction.}; that is the Word performing what appertains to the Word, and the flesh carrying out what appertains to the flesh. One of them sparkles with miracles, the other succumbs to injuries. And as the Word does not cease to be on an equality with His Father’s glory, so the flesh does not forego the nature of our race. For it must again and again be repeated that one and the same is truly Son of \textsc{God} and truly son of man. \textsc{God} in that “in the beginning was the Word, and the Word was with \textsc{God}, and the Word was \textsc{God\footnote{S. John i. l.}};” man in that “the Word became flesh and dwelt in us\footnote{Ibid. 14.}.” \textsc{God} in that “all things were made by Him\footnote{Ibid. 3, the Latin is \emph{per ipsum} (Gk. \textgreek{δι᾽ αὐτοῦ}) (through Him).}, and without Him was nothing made:” man in that “He was made of a woman, made under law\footnote{Gal. iv. 4.}.” The nativity of the flesh was the manifestation of human nature: the childbearing of a virgin is the proof of Divine power. The infancy of a babe is shown in the humbleness of its cradle\footnote{Viz., that it was laid “in a manger:” the Gk. version has \textgreek{σπαργάνων}, “swaddling clothes,” to represent \emph{cunarum} and this meaning is adopted by Bright [and Heurtley], S. Luke ii. 7.}: the greatness of the Most High is proclaimed by the angels’ voices\footnote{Ibid. 13.}. He whom Herod treacherously endeavours to destroy is like ourselves in our earliest stage\footnote{\emph{Similis est rudimentis hominum}.}: but He whom the Magi delight to worship on their knees is the \textsc{Lord} of all. So too when He came to the baptism of John, His forerunner, lest He should not be known through the veil of flesh which covered His Divinity, the Father’s voice thundering from the sky, said, “This is My beloved Son, in whom I am well pleased\footnote{S. Matt. iii. 17.}.” And thus Him whom the devil’s craftiness attacks as man, the ministries of angels serve as \textsc{God}. To be hungry and thirsty, to be weary, and to sleep, is clearly human: but to satisfy 5,000 men with five loaves, and to bestow on the woman of Samaria living water, droughts of which can secure the drinker from thirsting any more, to walk upon the surface of the sea with feet that do not sink, and to quell the risings of the waves by rebuking the winds, is, without any doubt, Divine. Just as therefore, to pass over many other instances, it is not part of the same nature to be moved to tears of pity for a dead friend, and when the stone that closed the four-days’ grave was removed, to raise that same friend to life with a voice of command: or, to hang on the cross, and turning day to night, to make all the elements tremble: or, to be pierced with nails, and yet open the gates of paradise to the robber’s faith: so it is not part of the same nature to say, “I and the Father are one,” and to say, “the Father is greater than I\footnote{S. John xiv. 28; x. 30: the reconciliation of this class of apparently contradictory statements is often undertaken by Leo [e.g. Sermon xxiii. 2 and lxxvii. 5; Ep. xxviii. 4 and lix. 3], and by other fathers (e.g. by Augustine \emph{de Fide et Symbolo,} 18).}.” For although in the \textsc{Lord} Jesus Christ \textsc{God} and man is one person, yet the source of the degradation, which is shared by both, is one, and the source of the glory, which is shared by both, is another. For His manhood, which is less than the Father, comes from our side: His Godhead, which is equal to the Father, comes from the Father.

V. \emph{Christ’s flesh is proved real from Scripture.}

Therefore in consequence of this unity of person which is to be understood in both natures\footnote{This is what theologians call \emph{communicatio idiomatum}, or in Gk. \textgreek{ἀντίδοσις}, the interchange of the properties of the two natures in Christ. The passage from the beginning of the chapter to “the Lord of glory” is somewhat freely adapted from S. Aug., c. Serm. Arian., cap. 8.}, we read of the Son of Man also descending from heaven, when the Son of \textsc{God} took flesh from the Virgin who bore Him. And again the Son of \textsc{God} is said to have been crucified and buried, although it was not actually in His Divinity whereby the Only-begotten is co-eternal and con-substantial with the Father, but in His weak human nature that He suffered these things. And so it is that in the Creed also we all confess that the Only-begotten Son of \textsc{God} was crucified and buried, according to that saying of the Apostle: “for if they had known, they would never have crucified the \textsc{Lord} of glory\footnote{1 Cor. ii. 8.}.” But when our \textsc{Lord} and Saviour Himself would instruct His disciples’ faith by His questionings, He said, “Whom do men say that I, the Son of Man, am?” And when they had put on record the various opinions of other people, He said, “But \emph{ye}, whom do ye say that I am?” Me, that is, who am the Son of Man, and whom ye see in the form of a slave, and in true flesh, whom do ye say that I am? Whereupon blessed Peter, whose divinely inspired confession was destined to profit all nations, said, “Thou art Christ, the Son of the living \textsc{God\footnote{S. Matt. xvi. 13–16.}}.” And not undeservedly was he pronounced blessed by the \textsc{Lord}, drawing from the chief corner-stone\footnote{\emph{A principali petra}. The Gk. version giving \textgreek{ἀπὸ τῆς πρωτοτύπου πέτρας}: others translate it “from the original (or archetypal) rock,” but it seems better to link the passage more closely with Eph. ii. 20; 1 Pet. ii. 6, \&c., although the Greek rendering is against this: see Serm. iv. chap. 2, where Leo is expounding the same favourite text. Bright’s note 64 is most useful in explaining the Leonine exposition. “Three elements,” he says, combine in the idea; (1) Christ Himself; (2) the faith in Christ; and (3) Peter considered as the chief of the Apostles and under Christ, the head of the Church.” Hence \emph{petra} is applied to each of these at different times.} the solidity of power which his name also expresses, he, who, through the revelation of the Father, confessed Him to be at once Christ and Son of \textsc{God}: because the receiving of the one of these without the other was of no avail to salvation, and it was equally perilous to have believed the \textsc{Lord} Jesus Christ to be either only \textsc{God} without man, or only man without \textsc{God}. But after the \textsc{Lord’s} resurrection (which, of course, was of His true body, because He was raised the same as He had died and been buried), what else was effected by the forty days’ delay than the cleansing of our faith’s purity from all darkness? For to that end He talked with His disciples, and dwelt and ate with them, He allowed Himself to be handled with diligent and curious touch by those who were affected by doubt, He entered when the doors were shut upon the Apostles, and by His breathing upon them gave them the Holy Spirit\footnote{S. John xx. 22.}, and bestowing on them the light of understanding, opened the secrets of the Holy Scriptures\footnote{S. Luke xxiv. 27.}. So again He showed the wound in His side, the marks of the nails, and all the signs of His quite recent suffering, saying, “See My hands and feet, that it is I. Handle Me and see that a spirit hath not flesh and bones, as ye see Me have\footnote{Ibid. 39.};” in order that the properties of His Divine and human nature might be acknowledged to remain still inseparable: and that we might know the Word not to be different from the flesh, in such a sense as also to confess that the one Son of \textsc{God} is both the Word and flesh\footnote{i.e. not to fall into the Charybdis of Nestorianism in avoiding the Scylla of Eutychianism.}. Of this mystery of the faith\footnote{\emph{Fidei sacramento}.} your opponent Eutyches must be reckoned to have but little sense if he has recognized our nature in the Only-begotten of \textsc{God} neither through the humiliation of His having to die, nor through the glory of His rising again. Nor has he any fear of the blessed apostle and evangelist John’s declaration when he says, “every spirit which confesses Jesus Christ to have come in the flesh, is of \textsc{God}: and every spirit which destroys Jesus is not of \textsc{God}, and this is Antichrist\footnote{John iv. 2, 3: the Lat. for “destroys” (or “dissolves,” Bright) is \emph{solvit} (so also in Lett. CXLIV. 3), which appears to be an exclusively Western reading: for Socrates, “the only Greek authority for \textgreek{λύει}” (the Gk. equivalent), according to Dr. Westcott, quotes no Gk. \textsc{mss.} as giving it, though he unhesitatingly makes use of that reading. The Gk. version here however, gives \textgreek{διαιρεῖν}, which simply begs the question (in Leo’s favour) as to the original meaning of the phrase \emph{solvere Jesum}, though on the face of it that is not at all necessarily obvious.}.” But what is “to destroy Jesus,” except to take away the human nature from Him, and to render void the mystery, by which alone we were saved, by the most barefaced fictions. The truth is that being in darkness about the nature of Christ’s body, he must also be befooled by the same blindness in the matter of His sufferings. For if he does not think the cross of the \textsc{Lord} fictitious, and does not doubt that the punishment He underwent to save the world is likewise true, let him acknowledge the flesh of Him whose death he already believes: and let him not disbelieve Him man with a body like ours, since he acknowledges Him to have been able to suffer: seeing that the denial of His true flesh is also the denial of His bodily suffering. If therefore he receives the Christian faith, and does not turn away his ears from the preaching of the Gospel: let him see what was the nature that hung pierced with nails on the wooden cross, and, when the side of the Crucified was opened by the soldier’s spear, let him understand whence it was that blood and water flowed, that the Church of \textsc{God} might be watered from the font and from the cup\footnote{\emph{Et lavacro rigaretur et poculo}: that is by the two great “generally necessary” sacraments of which he takes the water and the blood “from His riven side which flowed” to be a symbol.}. Let him hear also the blessed Apostle Peter, proclaiming that the sanctification of the Spirit takes place through the sprinkling of Christ’s blood\footnote{This refers to 1 Pet. i. 2 (\emph{q.v.}).}. And let him not read cursorily the same Apostle’s words when he says, “Knowing that not with corruptible things, such as silver and gold, have ye been redeemed from your vain manner of life which is part of your fathers’ tradition, but with the precious blood of Jesus Christ as of a lamb without spot and blemish\footnote{1 Pet. i. 18.}.” Let him not resist too the witness of the blessed Apostle John, who says: “and the blood of Jesus the Son of \textsc{God} cleanseth us from all sin\footnote{1 John i. 7.}.” And again: “this is the victory which overcometh the world, our faith.” And “who is He that overcometh the world save He that believeth that Jesus is the Son of \textsc{God}. This is He that came by water and blood, Jesus Christ: not by water only, but by water and blood. And it is the Spirit that testifieth, because the Spirit is the truth\footnote{Some of the \textsc{mss.} here give \emph{Christus} for \emph{Spiritus} (the reading adopted also by the Vulgate): in this case you must translate \emph{that Christ} is the Truth instead of \emph{because of the Spirit, \&c}.: but see Westcott’s note \emph{in loc}.}, because there are three that bear witness, the Spirit, the water and the blood, and the three are one\footnote{1 John v. 4–8. The absence of the verse on the “Heavenly witnesses” (distinctly a western insertion) is to be noticed. On Leo’s interpretation of this mysterious passage Canon Bright’s note 168 should be consulted.}.” The Spirit, that is, of sanctification, and the blood of redemption, and the water of baptism: because the three are one, and remain undivided, and none of them is separated from this connection; because the catholic Church lives and progresses by this faith, so that in Christ Jesus neither the manhood without the true Godhead nor the Godhead without the true manhood is believed in.

VI. \emph{The wrong and mischievous confession of Eutyches. The terms on which he may be restored to Communion. The sending of deputies to the east.}

But when during your cross-examination Eutyches replied and said, “I confess that our \textsc{Lord} had two natures before the union but after the union I confess but one\footnote{This was the only compromise of his views which Eutyches could be brought to make at the synod of Constantinople. Though it was rejected, and did not hinder his condemnation, it was never met with a direct, categorical refutation.},” I am surprised that so absurd and mistaken a statement of his should not have been criticised and rebuked by his judges, and that an utterance which reaches the height of stupidity and blasphemy should be allowed to pass as if nothing offensive had been heard: for the impiety of saying that the Son of \textsc{God} was of two natures before His incarnation is only equalled by the iniquity of asserting that there was but one nature in Him after “the Word became flesh.” And to the end that Eutyches may not think this a right or defensible opinion because it was not contradicted by any expression of yourselves, we warn you beloved brother, to take anxious care that if ever through the inspiration of \textsc{God’s} mercy the case is brought to a satisfactory conclusion, his ignorant mind be purged from this pernicious idea as well as others. He was, indeed, just beginning to beat a retreat from his erroneous conviction, as the order of proceedings shows\footnote{\emph{Gestorum ordo}, as before, in chap. 1. A report of the proceedings had accompanied Flavian’s letter.}, in so far as when hemmed in by your remonstrances he agreed to say what he had not said before and to acquiesce in that belief to which before he had been opposed. However, when he refused to give his consent to the anathematizing of his blasphemous dogma, you understood, brother\footnote{\emph{Fraternitas vestra}: or, as the Gk. version apparently took it, “you and the rest of the brethren” (\textgreek{ἡ ὑμῶν ἀδελφότης}).}, that he abode by his treachery and deserved to receive a verdict of condemnation. And yet, if he grieves over it faithfully and to good purpose, and, late though it be, acknowledges how rightly the bishops’ authority has been set in motion; or if with his own mouth and hand in your presence he recants his wrong opinions, no mercy that is shown to him when penitent can be found fault with\footnote{It will be remembered that he had been degraded from the priesthood and deprived of his monastery, as well as excommunicated: he might be reinstated in all these privileges, the mercifulness of Leo hints, if he recant his errors.}: because our \textsc{Lord}, that true and “good shepherd” who laid down His life for His sheep\footnote{S. John x. 11 and 15.} and who came to save not lose men’s souls\footnote{S. Luke ix. 50.}, wishes us to imitate His kindness\footnote{\emph{Pietatis}, a beautiful word, expressing now the Father’s pitying protection, now the children’s loyal affection, and here the Elder Brother’s love for the younger and weaker. Cf. n. I. on chap. iii.}; in order that while justice constrains us when we sin, mercy may prevent our rejection when we have returned. For then at last is the true Faith most profitably defended when a false belief is condemned even by the supporters of it.

Now for the loyal and faithful execution of the whole matter, we have appointed to represent us our brothers Julius\footnote{Bishop of Puteoli.} Bishop and Renatus\footnote{Died at Delos on the way. The words “of the title of S. Clement” are of doubtful authenticity, and not found in the Gk. version. The parish churches of Rome seem to have been called \emph{tituli} at their first founding about the beginning of the 4th cent. \textsc{a.d.} Cf. our Eng. term “title,” and refer to Bingham, Bk. viii. § 1.} priest [of the Title of S. Clement], as well as my son Hilary\footnote{Afterwards Leo’s successor in the see of Rome, 461–8.}, deacon. And with them we have associated Dulcitius our notary, whose faith is well approved: being sure that the Divine help will be given us, so that he who had erred may be saved when the wrongness of his view has been condemned. \textsc{God} keep you safe, beloved brother.

The 13 June, 449, in the consulship of the most illustrious Asturius and Protogenes.
